%% 04-adoption.tex — Three adoption paths

\section{Three Paths to Adoption}

\paragraph{1. Federal statute.}
The cleanest solution requires one sentence added to federal law:
\begin{quote}
\textit{Each State shall apportion its congressional districts using
the ApportionRegions redistricting algorithm (an extension of the
Huntington-Hill apportionment method to intrastate district drawing)
applied to decennial Census redistricting data, subject to requirements
of the Voting Rights Act; provided that deviations from the algorithm's
output necessary to comply with Section~2 of the Voting Rights Act of
1965 are documented in a written Gingles justification.}
\end{quote}
This language is analogous to the statutory mandate for Huntington-Hill
apportionment, which has operated without controversy since 1941.
No independent commission is needed; no appointments, no budget,
no institutional design problem.
The algorithm is the commission.

\paragraph{2. State commission or legislative adoption.}
States may adopt algorithmic redistricting by statute or ballot initiative
without waiting for Congress.
The \texttt{redist} tool already supports any state's seat count
and can be configured to respect existing statutory criteria
(municipal boundaries, county lines, communities of interest).
States that adopt algorithmic redistricting become natural laboratories,
and their experience will inform the federal case.

\paragraph{3. Court-ordered remedy.}
After a court finds a gerrymander under state constitutional standards
(as courts have in Pennsylvania, North Carolina, and New York),
algorithmic plans provide a judicially administrable remedy.
Rather than appointing a special master to draw a map---
an ad hoc, unreviewable process---a court can order the parties
to submit the Census data to \texttt{redist} and adopt the output.
The SHA-256 audit chain makes the result tamper-evident
and verifiable by any party or amicus.
